\documentclass[12pt]{article}

%\usepackage[cm]{fullpage}
\setlength{\topmargin}{-1in}
\setlength{\oddsidemargin}{-0.25in}
\setlength{\evensidemargin}{-0.25in}
\setlength{\textwidth}{7in}
\setlength{\textheight}{10in}

\usepackage{array}
\usepackage{enumitem}
\usepackage{amssymb}
\usepackage{amsmath}
\usepackage{epstopdf}
\usepackage[T1]{fontenc}
\usepackage{graphics}
\usepackage{graphicx}
\usepackage{hyperref}
\usepackage{latexsym}
\usepackage{mathrsfs}
\usepackage{multicol}
\usepackage{setspace}
\usepackage{subfigure}
\usepackage{time}
\usepackage{times}

\newcommand{\To}{\Rightarrow}
\newcommand{\del}{\nabla}
\newcommand{\R}{\mathbb{R}}
\newcommand{\C}{\mathcal{C}}
\newcommand{\N}{\mathcal{N}}
\newcommand{\A}{\mathcal{A}}
\renewcommand{\O}{\mathcal{O}}
% \renewcommand{\N}{\mathbb{N}}
\newcommand{\ben}{\begin{enumerate}}
\newcommand{\een}{\end{enumerate}}
\newcommand{\eps}{\varepsilon}
\renewcommand{\epsilon}{\varepsilon}
\newcommand{\p}{\partial}
\newcommand{\<}{\langle}
\renewcommand{\>}{\rangle}
\renewcommand{\subset}{\subseteq}
\newcommand{\cont}{\Rightarrow \Leftarrow}
\newcommand{\back}{\backslash}
\newcommand{\norm}[1]{\|{#1}\|}
\newcommand{\abs}[1]{|{#1}|}
\newcommand{\ip}[1]{\langle{#1}\rangle}
\newcommand{\m}[1]{{\mathbf{#1}}}
\newcommand{\sts}{\m{S}^\textrm{T}\m{S}}
\newcommand{\rank}{\textrm{rank}}
\newcommand{\X}{\mathcal{X}}
\newcommand{\Y}{\mathcal{Y}}
\def\ds{\displaystyle}
\def\d{\partial}
\renewcommand{\arraystretch}{1.5}
\renewcommand{\L}{\mathscr{L}}

%\includegraphics[scale = 0.65, trim = 1in 3.25in 1in 3in]{blank_axes.pdf}


\pagestyle{empty}
\begin{document}
\noindent\rule{7in}{0.4pt}\\
\begin{center}
{\sc Differential Equations\\
Linearity}
\end{center}
\noindent\rule{7in}{0.4pt}
%%%%%%%%%%%%%%%%%%%%%%%%%%%%%%%%%%%%%%%%%%%%%%%%
\section*{Classifying DEs:  linear vs. nonlinear}
\doublespacing
\noindent  Another classification that will be used to determine which technique we can use to solve a
DE is whether a DE is {\bf linear or nonlinear}.  The technical definition of linearity is not especially
easy to un-pack.Here it is, but then we need to discuss what it means.
	\begin{quote}
	A DE %with independent variable $x$ and dependent variable $y$ 
	is said to be {\it linear} if it can be written as \[ a_n(x)y^{(n)} + a_{n-1}y^{(n-1)} + \cdots + a_2(x)y'' + a_1(x)y' + a_0(x)y = f(x). \] A DE is said to be {\it nonlinear} if it is not linear.
	\end{quote}
\noindent  This definition is a good theoretical/mathematical definition, but what are the important
features we can tease out of the definition?  In this case, it might be easier to describe what would
be true of a DE that does not fit in the form described in the definition of linearity.\\

\noindent  A DE is {\bf nonlinear} if:
	\begin{itemize}
	\item \mbox{}\\ \vfill
		%The dependent variable or a derivative is raised to a power other than 1.
		%For example, the DE $y'' + y^2 = x^3$
		%is nonlinear because the dependent variable, $y$, is raised to the second power.
		
	\item   \mbox{}\\ \vfill
		%The dependent variable or a derivative is inside a function.  In the DE $y'' + \ln(y') = e^x$,
		%the derivative $y'$ is inside a logarithmic function, so this DE is nonlinear.
	\item   \mbox{}\\ \vfill
		%The dependent variable and/or any of its derivatives are attached to each other via multiplication.
		%For example,the DE $y'\cdot y'' + e^x y = y',$ is nonlinear because the first term is a product 
		%of $y'$ and $y''$.
	\end{itemize}
	


\end{document}